\subsection{P033 --- Enabling scalable Model Predictive Control design for building HVAC systems using semantic data modelling}
\label{paper:P033}
\textbf{Authors:} Lu Wan, Ferdinand Rossa, Torsten Welfonder, Ekaterina Petrova, Pieter Pauwels\par
\textbf{Major Category:} BIM and Semantic Information\par
\textbf{Secondary Category:} BMS, BAS, BACS and Building Automation, Building Energy and ZEB\par
\textbf{Keywords:} Model Predictive Control, Semantic data modelling, Ontologies, Building Information Modelling, HVAC, RDF, System identification\par
\paragraph{主な主張}
既存ontologyを組み合わせたshareable semantic data modelとBIM-to-RDF toolingにより, MPC configurationに必要なmetadataを統合し, semantic graphからcontrol-oriented RC modelのsystem identificationを自動化できることを実建物で示した. \cite{Wan:2025SemanticMPC}
\paragraph{新規性}
semantic modellingをBMS pointのinteroperabilityに留めず, building physics, geometry, HVAC topology, sensor/actuator, MPC algorithm metadataまで統合し, control model configurationとparameterizationへ直接接続した点に新規性がある.
\paragraph{位置付け}
semantic graphからsystem identificationとcontrol model configurationへ接続する代表例であり, semanticsをprediction/controlの前処理基盤からmodel-generation layerへ拡張する研究として位置付ける.
